% !PRINTED 118
\weberpage{118}
%<<CONT
ein Element $b$ heraus, so ist der Complex $Pb$ mit $P$ völlig
identisch. Es können sich also die beiden Reihen
\[
A = a_0,\ a_1,\ a_2,\ \ldots,\ a_{n-1}
\]
\[
Ab = a_0b,\ a_1b,\ a_2b,\ \ldots,\ a_{n-1}b
\]
nur durch die Anordnung von einander unterscheiden. Der
Uebergang von $A$ zu $Ab$ ist also eine Permutation von $n$ Ziffern
$0, 1, 2 \ldots n - 1$, und jedem Elemente $b$ von $P$ entspricht eine
solche Permutation, die wir mit $\pi_b$ bezeichnen wollen. Zwei verschiedenen
Elementen $b, c$ entsprechen immer zwei verschiedene
Permutationen $\pi_b, \pi_c$, und dem Einheitselemente entspricht die
identische Permutation. Setzen wir
\[
\pi_b = (A,\ Ab), \quad \pi_c = (A,\ Ac),
\]
so können wir $\pi_c$ auch mit $(Ab, Abc)$ bezeichnen, und daraus
ergiebt sich
\weq{2}{\pi_b \pi_c = \pi_{bc},}
d.\,h. die Permutationen $\pi$ bilden eine mit $P$ isomorphe Permutationsgruppe. Diese Permutationsgruppe ist transitiv, da z.\,B.
das Element $a_0$ in jedes beliebige andere Element von $P$ übergehen kann. Also haben wir den Satz:

1.\ \gesperrt{Jede Gruppe vom Grade $n$ ist isomorph mit einer
transitiven Permutationsgruppe von $n$ Ziffern.}

Es giebt natürlich noch andere mit einer gegebenen Gruppe
isomorphe Permutationsgruppen, und es wäre von besonderem
Interesse, eine solche Gruppe mit möglichst geringer Ziffernzahl
zu bilden. Diese Aufgabe kann bis jetzt nicht allgemein gelöst
werden. Wir müssen uns hier mit wenigen Sätzen begnügen.

Wir nehmen an, es sei $R$ irgend ein Divisor von $P$ vom
Index $j$, und bezeichnen die Elemente von $R$ mit $c$, setzen ferner,
indem wir $P$ in ein System von Nebengruppen zerlegen,
\weq{3}{P = R + Rb_1 + Rb_2 + \cdots + Rb_{j-1}.}
Ist dann $a$ irgend ein Element von $P$, so werden die beiden
Systeme
\weq{4}{\begin{aligned}
B &= R,\ Rb_1,\ Rb_2,\ \ldots,\ Rb_{j-1}\\
Ba &= Ra,\ Rb_1a,\ Rb_2a,\ \ldots,\ Rb_{j-1}a,
\end{aligned}}
von der Reihenfolge abgesehen, übereinstimmen, und es entspricht also jedem Element $a$ eine bestimmte Permutation $\pi_a$ der
$j$ Ziffern $0, 1 \ldots j - 1$, die wir mit
\weq{5}{\pi_a = (B,\ Ba)}
%CONT>>
